\begin{lemma}
\label{editorial-source-lemma-powers-field}
\begin{reference}
\cite[Lemma 3.1.6]{Alper-adequate}
\end{reference}
Let $k'/k$ be a field extension. The following are equivalent
\begin{enumerate}
\item for each $x \in k'$ there exists an $n > 0$ such that $x^n \in k$, and
\item $k' = k$ or $k$ and $k'$ have characteristic $p > 0$ and
either $k'/k$ is a purely inseparable extension or
$k$ and $k'$ are algebraic extensions of $\mathbf{F}_p$.
\end{enumerate}
\end{lemma}

\begin{proof}
Each case of (2) implies (1). For $k'=k$
take $n=1$. For a purely inseparable
extension in characteristic $p>0$, each
original element has a $p$-power in $k$.
If both fields are algebraic over
$\mathbf F_p$, a nonzero $x\in k'$
lies in the finite field $\mathbf F_p(x)$,
whose multiplicative group is finite;
some positive power of $x$ is $1\in k$.
For $x=0$ take $n=1$.

Assume (1) and retain the original extension
$k'/k$. Each $x\in k'$ is algebraic over
$k$, being a root of $T^n-x^n\in k[T]$.
Let $E$ be its maximal separable
subextension from Fields, Lemma
\ref{fields-lemma-separable-first}.
In positive characteristic, if $E=k$
then $k'/k$ is purely inseparable and
we have finished. In characteristic
zero $E=k'$, so $E=k$ is precisely the
identity case. Otherwise choose
$x\in E\setminus k$ and keep the finite
separable intermediate field $L=k(x)$.
Property (1) is inherited by $L/k$.
It is enough to show that $k$ has
positive characteristic and is algebraic
over its prime field: the original
$k'/k$ is already algebraic, and then
$k'$ is algebraic over the same prime
field by transitivity. This proves the
reduction for infinite $k'/k$ as well.

First let $z$ be any element of $L\setminus k$.
There are $k$-embeddings $\sigma,\tau$ of
$k(z)$ into an algebraic closure $\overline k$
with $A=\sigma(z)\neq B=\tau(z)$, because
its separable minimal polynomial has
degree at least two. This is the
embedding count of Fields, Lemma
\ref{fields-lemma-count-embeddings}.
Both $z$ and $z+1$ are nonzero, since
neither belongs to $k$. Choose the
original positive integers $n,m$ with
$z^n,(z+1)^m\in k$. Their images give
$A^n=B^n$ and $(A+1)^m=(B+1)^m$.
Thus, with all denominators nonzero,
$$
\zeta=A/B,\qquad \zeta'=(A+1)/(B+1),
\qquad \zeta^n=1,\quad(\zeta')^m=1.
$$
One has $\zeta'\neq1$ because $A\neq B$.
The full identity
$$
\zeta'(B+1)=A+1=\zeta B+1
$$
gives $B(\zeta'-\zeta)=1-\zeta'$.
Since the right side is nonzero,
$\zeta'-\zeta\neq0$, and hence
$$
B=\frac{1-\zeta'}{\zeta'-\zeta}.
$$
Both roots of unity are algebraic over
the prime field, so this formula makes
$B$ algebraic over that field. The
embedding $\tau$ fixes the prime field;
pulling back an annihilating polynomial
through the injective $\tau$ proves that
the original $z$ is algebraic over it.
This applies to every $z\in L\setminus k$.
For any $a\in k$, both $x$ and $x+a$
lie outside $k$ and are algebraic over
the prime field, so their difference
$a=(x+a)-x$ is algebraic. It follows
that $k$, and then the original $k'$,
are algebraic over the prime field.

It remains to exclude characteristic
zero. Use the fixed original $x\in L\setminus k$
and two embeddings as above, with
$A=\sigma(x)$, $B=\tau(x)$, and the
same $\zeta,\zeta'$. Both $A=\zeta B$
and $B=(1-\zeta')/(\zeta'-\zeta)$
belong to the number field
$F=\mathbf Q(\zeta,\zeta')$.
This field contains only finitely many
roots of unity, for the following exact
reason. Put $d=[F:\mathbf Q]$. The monic
minimal polynomial of a root of unity
in $F$ has degree $e\leq d$ and integer
coefficients: it divides some $T^N-1$
in $\mathbf Q[T]$, and the primitive
Gauss lemma gives a factor in $\mathbf Z[T]$
whose leading coefficient divides $1$.
Every complex root $\xi_i$ of that
minimal polynomial is also an $N$th
root of unity and has absolute value
$1$. Its coefficient of $T^{e-j}$ is
the full signed sum
$$
(-1)^j\sum_{1\leq i_1<\cdots<i_j\leq e}
\xi_{i_1}\cdots\xi_{i_j},
$$
so its absolute value is at most
$\binom ej$. There are finitely many
integer coefficient lists with these
bounds and $e\leq d$, and each resulting
nonzero polynomial has finitely many
roots. This proves the required
finiteness inside $F$.

For every original rational $y$, one
has $x+y\notin k$ and $B+y\neq0$.
Property (1) gives a positive exponent
for $x+y$, so
$$
\zeta_y=\frac{A+y}{B+y}
=\frac{\zeta B+y}{B+y}\in F
$$
is a root of unity, and $\zeta_y\neq1$.
Infinitely many rational $y$ and only
finitely many roots of unity in $F$
give distinct $y,y'$ with
$\zeta_y=\zeta_{y'}$. Subtracting the
two original identities gives
$$
y-y'=(A+y)-(A+y')
=\zeta_y\big((B+y)-(B+y')\big)
=\zeta_y(y-y').
$$
Since $y-y'\neq0$, this forces
$\zeta_y=1$, a contradiction. Thus the
nontrivial separable case has positive
characteristic, and the already proved
algebraicity gives exactly (2).

The argument also determines when a
single positive exponent can work for
all elements. Such an exponent exists
if and only if $k'=k$, or the fields
have characteristic $p>0$ and either
$k'/k$ is purely inseparable of bounded
exponent, or $k'$ is a finite field.
The converse is immediate with exponent
$1$, a common $p$-power, or $|k'|-1$,
respectively. For necessity retain a
single original exponent $N$ for all
elements. If a nontrivial separable
intermediate field exists, choose
$A=\sigma(x)\neq B=\tau(x)$ as above.
For every $a\in k$, the quotient
$(A+a)/(B+a)$ is an $N$th root of unity.
These quotients are distinct for
distinct $a$: cross multiplication
gives the full identity
$$
(A+a)(B+b)-(A+b)(B+a)=(A-B)(b-a).
$$
Their denominators are nonzero because
$x\notin k$. Hence $|k|\leq N$, as a
nonzero degree-$N$ polynomial has at
most $N$ roots. The field $k$ is finite.
Every element of $k'$ is then among
the roots of the finitely many original
polynomials $T^N-c$, $c\in k$. For
$c=0$ the only root is zero; for
each of the $|k|-1$ nonzero values
there are at most $N$ roots. Thus
$|k'|\leq1+N(|k|-1)$, and $k'$ is finite.

In the remaining nonidentity case
$k'/k$ is purely inseparable in
characteristic $p$. Write the original
$N=p^e m$ with $p\nmid m$. For every
nonzero $x\in k'$, put $u=x^{p^e}$.
It is purely inseparable over $k$ and
satisfies $u^m=x^N\in k$. The polynomial
$T^m-x^N$ has nonzero derivative at
each root, since $m$ and each root
are nonzero. Hence $u$ is also
separable over $k$, and belongs to
$k$. The zero element satisfies the
same conclusion directly. Therefore
$(k')^{p^e}\subseteq k$, giving a
uniform inseparability exponent from
the full original integer $N$.

For finite original fields the fixed-exponent
condition has an exact form. Put $q=|k|$
and $q'=|k'|$. The multiplicative group
of a finite field is cyclic, as follows
directly from the root bound. For a
finite subgroup $G$ of a field's
multiplicative group, let $h$ be the
least common multiple of its element
orders. For every prime $\ell$ dividing
$h$, choose an element whose order
has the largest $\ell$-power $\ell^{a_\ell}$
occurring in $h$, and raise that element
to the prime-to-$\ell$ part of its
order. The result has order exactly
$\ell^{a_\ell}$. The product of these
elements has order $h$: if its $m$th
power is $1$, raising to
$h/\ell^{a_\ell}$ shows
$\ell^{a_\ell}\mid m$ for each $\ell$,
since $h/\ell^{a_\ell}$ is coprime
to $\ell$. Conversely its $h$th power
is $1$. All elements of $G$ are
roots of $T^h-1$, so $|G|\leq h$;
the element of order $h$ gives
$|G|\geq h$. Thus it generates $G$.
For the trivial group the empty
product has order $1$, giving the
same conclusion.

Choose a generator $g$ of the original
$(k')^\times$, of order $q'-1$.
Inside $k'$, the nonzero roots of
$T^{q-1}-1$ are exactly $k^\times$:
all its $q-1$ elements are roots,
and a degree-$q-1$ polynomial has
no additional ones. Therefore
\[

\begin{split}
(\forall x\in k')\ x^N\in k
&\quad\Longleftrightarrow\quad g^N\in k^\times\\
&\quad\Longleftrightarrow\quad g^{N(q-1)}=1\\
&\quad\Longleftrightarrow\quad (q'-1)\mid N(q-1).
\end{split}

\]
The zero element satisfies the condition
for every positive $N$. If $d=[k':k]$,
then $q'=q^d$, by counting the coordinate
vectors of a $k$-basis. Retaining the
full factorization
$$
q'-1=q^d-1=(q-1)(1+q+\cdots+q^{d-1}),
$$
the displayed divisibility is equivalently
$(1+q+\cdots+q^{d-1})\mid N$. The least
positive common exponent is consequently
the full sum $1+q+\cdots+q^{d-1}$.
This also proves the sharp cardinality
bound above: equality holds for this
least exponent, with every contribution
from the original zero and nonzero
elements retained.
\end{proof}